\documentclass{article} % 12pt is optional
\usepackage[utf8]{inputenc}
\usepackage[a4paper,margin=2cm]{geometry} 
\usepackage{amssymb, amsmath, amsthm}
\usepackage[most]{tcolorbox}
\usepackage[dvipsnames,svgnames,x11names]{xcolor}
\usepackage{fullpage}
\usepackage[catalan]{babel}
\usepackage[colorlinks=true, allcolors=black]{hyperref}
\usepackage{mathtools}
\newcommand{\R}{\mathbb{R}}
\newcommand{\Z}{\mathbb{Z}}
\usepackage{physics}    
\usepackage{siunitx}     
\usepackage{enumitem}
\usepackage{microtype}   
\usepackage{tikz}
\usepackage{pgfplots}
\pgfplotsset{compat=1.18} 
\graphicspath{ {./images/} }
\usepackage{booktabs}
\setlength{\heavyrulewidth}{0.4pt}
\linespread{1.25}
\setlength{\lightrulewidth}{0.4pt}
\setlength{\cmidrulewidth}{0.4pt}
\usepackage{chemfig}
\usepackage[version=4]{mhchem}
\usepackage[table]{xcolor}
\title{\textbf{ESQUEMA MATEMÀTIQUES}}
\author{\textit{Gerard Cosp Lores}}
\usepackage{graphicx}
\usepgfplotslibrary{fillbetween}
\usepackage{float}
\date{}

\begin{document}

\maketitle

\begin{abstract}
Aquest document està destinat a millorar els errors que es cometen en els exàmens de matemàtiques per \textbf{poder-ho fer millor}\footnote{Aquest document s'anirà actualitzant cada tema}.
\end{abstract}

\section{T1-T2 Matrius i determinants}
\subsection{Estudiar \textbf{Rang}(A) amb paràmetre}

\begin{tcolorbox}[
  title=Mètode del pivot,
  colback=blue!5!white,
  colframe=blue!70!black,
  coltitle=white,
  colbacktitle=blue!70!black,
  fonttitle=\bfseries,
  arc=6pt,
  boxrule=0.8pt
]

Primer escollim un menor d'ordre dos 
$\begin{vmatrix}
a & b \\
c & d
\end{vmatrix} \neq 0$ \textit{(millor si no te cap paràmetre)} i després resolem 1 \textbf{d'ordre 3} i l'altre ho fem \textbf{a partir dels resultats de l'anterior.}


\end{tcolorbox}

\begin{align}
A = \begin{vmatrix}
1 & 2 & a\\
a-1 & 2a & 0 \\
1 & -4 & -a
\end{vmatrix}
\end{align}

\begin{align}
\begin{vmatrix}
1 & 2 \\
a-1 & 2a
\end{vmatrix} = 9 \neq 0 \therefore Rang(A) \geq 2
\end{align}

\begin{align}
\begin{vmatrix}
1 & 2 & a\\
a-1 & 2a & 0 \\
1 & -4 & -a
\end{vmatrix} = (-2a^2 -4a^2+4a)-(2a^2 -2a^2 +2a)
\end{align}

\begin{align}
-6a^2+2a = 0 \therefore a = 0, a = \frac{1}{3}
\end{align}

Estudiariem els dos cassos si la matriu fos més gran, però hem obtingut que:

\begin{center}
\begin{tabular}{|c|c|}
\hline
\textbf{Valor de "$a$"} & \textbf{Rang de la matriu} \\ \hline
Si a = 0 i a = $\frac{1}{3}$ & Rang(A) = 2 \\ \hline
Si a $\neq$ 0 i a $\neq \frac{1}{3}$ & Rang(A) = 3 \\ \hline
\end{tabular}
\end{center}

\subsection{Estudiar inversa amb paràmetre}

\begin{tcolorbox}[
  title=Inversa amb determinant,
  colback=TealBlue!5!white,
  colframe=TealBlue!70!black,
  coltitle=white,
  colbacktitle=TealBlue!70!black,
  fonttitle=\bfseries,
  arc=6pt,
  boxrule=0.8pt
]

$\forall$ Matrius \textbf{$\exists$ inversa} si, i només si
$\begin{vmatrix}
a & b \\
c & d
\end{vmatrix} \neq 0$, de manera que el rang és el màxim, el que farem serà \textbf{igualar el determinant $= 0$} $. {\textbf{\textcolor{red}{Comprovar que és inversa de A i dona Id.}}}$ Podem fer la inversa de diverses maneres:

\begin{itemize}
    \item $A\cdot A^{-1} = A^{-1}\cdot A = Id$
    \item $A^{-1} = \frac{1}{\abs{A}} Adj(A^t)$
    \item Fem gauss amb $\begin{pmatrix}
a & b & c \mid & 1 & 0 & 0 \\
d & e & f \mid & 0 & 1 & 0 \\
g & h & i \mid & 0 & 0 & 1 \\
\end{pmatrix}$
\end{itemize}

\end{tcolorbox}

\subsection{Resoldre equacions matricials}

\begin{tcolorbox}[
  title=Equacions matricials,
  colback=MidnightBlue!5!white,
  colframe=MidnightBlue!70!black,
  coltitle=white,
  colbacktitle=MidnightBlue!70!black,
  fonttitle=\bfseries,
  arc=6pt,
  boxrule=0.8pt
]

Una equació matricial $\longrightarrow$ $X \cdot A = Id$ \textbf{es resol com una equació normal}, però no podem dividir per tant \textit{multipliquem} per $A^{-1}$ (que la podem fer per \textbf{definició, formula o Gauss}).

\end{tcolorbox}

\subsection{Altres tipus de matrius}

\begin{tcolorbox}[
  title=Tipus de matrius,
  colback=NavyBlue!5!white,
  colframe=NavyBlue!70!black,
  coltitle=white,
  colbacktitle=NavyBlue!70!black,
  fonttitle=\bfseries,
  arc=6pt,
  boxrule=0.8pt
]

\begin{itemize}
  \item Matriu Diagonal Superior $\begin{pmatrix}
    1 & 0 & 0\\
    3 & 2 & 0 \\
    1 & -4 & -1
  \end{pmatrix}$
  \item Matriu Diagonal Inferior $\begin{pmatrix}
    1 & 2 & 4\\
    0 & 2 & 0 \\
    0 & 0 & -1
  \end{pmatrix}$
  \item Matriu Identitat $\begin{pmatrix}
    1 & 0 & 0\\
    0 & 1 & 0 \\
    0 & 0 & 1
  \end{pmatrix}$
  \item Matriu Nula $\begin{pmatrix}
    0 & 0 & 0\\
    0 & 0 & 0 \\
    0 & 0 & 0
  \end{pmatrix}$
\end{itemize}
\end{tcolorbox}

\section{T3 Sistemes d'equacions}

\subsection{Discuteix el tipus de sistema i resol el sistema}

\begin{tcolorbox}[
  title=Teorema Rouché-Frobenius,
  colback=Salmon!5!white,
  colframe=Salmon!70!black,
  coltitle=white,
  colbacktitle=Salmon!70!black,
  fonttitle=\bfseries,
  arc=6pt,
  boxrule=0.8pt
]

El tipus de sistema \textit{SDC, SCI, SI}, \textbf{depen del Rang} de la matriu del sistema $\therefore$

\begin{itemize}
\item Si Rang($A$) = Rang($A^*$) = n(Incògnites)
$\Rightarrow$ SCD
\item Si Rang($A$) = Rang($A^*$) $<$ n(Incògnites)
$\Rightarrow$ SCI
\item Si Rang($A$) $\neq$ Rang($A^*$)
$\Rightarrow$ SI
\end{itemize}

\end{tcolorbox}

\noindent
1) Expressar el sistema en forma matricial\\
2) Estudiar el rang d'A a partir del pivot\\
3) Estudiar el rang de l'ampliada amb els valors del anterior\\
4) Explicar per quins valors es quin sistema i resoldre, (Cramer, Gauss /Paràmetres/)

\subsection{Resol un problema escrit}

\begin{tcolorbox}[
  title=Problemes escrits,
  colback=OrangeRed!5!white,
  colframe=OrangeRed!70!black,
  coltitle=white,
  colbacktitle=OrangeRed!70!black,
  fonttitle=\bfseries,
  arc=6pt,
  boxrule=0.8pt
]

Important \textbf{plantejar les variables} \textit{(+0.5 pts),} i \textbf{redactar un text final} en f(x) amb l'enunciat, intentar que els coeficients siguin $\mathbb{Z}$ per poder fer Gauss millor. \textit{Sobretot intentar, tot suma}.

\end{tcolorbox}

\section{T4 Limits i continuïtat}

\subsection{Resol els límits següents}

\begin{tcolorbox}[
  title=Resoldre limits,
  colback=SeaGreen!5!white,
  colframe=SeaGreen!70!black,
  coltitle=white,
  colbacktitle=SeaGreen!70!black,
  fonttitle=\bfseries,
  arc=6pt,
  boxrule=0.8pt
]

Sempre comprovar amb la calculadora, \textbf{apropant el valor de les x al limit i veient el comportament}. $\lim_{x \to \infty} f(x) \approx f(\infty)$, i enrecordar-se de que $1^{\infty} \approx e$.

\end{tcolorbox}

\subsection{Estudia la continuïtat de:}

\begin{tcolorbox}[
  title=Continuïtat,
  colback=ForestGreen!5!white,
  colframe=ForestGreen!70!black,
  coltitle=white,
  colbacktitle=ForestGreen!70!black,
  fonttitle=\bfseries,
  arc=6pt,
  boxrule=0.8pt
]

És important dir \textit{La funció es continua en tot el seu Domf, menys en els punts:} I \textbf{estudiem els punts crítics}, pot ser:

\begin{itemize}
\item Discontinuïtat evitable \textbf{(punt blanc)}
\item Discontinuïtat de salt infinit \textbf{(assimptota)}
\item Discontinuïtat de salt finit
\end{itemize}

\begin{center}
\begin{tikzpicture}
  \begin{axis}[
      width=3cm,
      height=2cm,
      scale only axis,
      axis lines=middle,              
      axis on top,                   
      axis line style={thick, black},
      xlabel={$x$}, ylabel={$f(x)$},  
      tick label style={font=\tiny},
      label style={font=\tiny},
      grid=both,                        
      domain=-3:3, samples=200,       
      ymin=-10, ymax=10,                
      xtick={-3,-2,-1,0,1,2,3},        
      ytick={-10,-5,0,5,10}
    ]
    \addplot[LimeGreen, line width=2pt, domain=-3:-0.1] {1/x};
    \addplot[LimeGreen, line width=2pt, domain=0.1:1] {1/x};
    \addplot[LimeGreen, line width=2pt, domain=1:3] {-x};
  \end{axis}
\end{tikzpicture}
\end{center}
\end{tcolorbox}

\subsection{Per a quins valors de a la funció es continua}

\begin{tcolorbox}[
  title=Límits amb paràmetres,
  colback=LimeGreen!5!white,
  colframe=LimeGreen!70!black,
  coltitle=white,
  colbacktitle=LimeGreen!70!black,
  fonttitle=\bfseries,
  arc=6pt,
  boxrule=0.8pt
]

Plantegem tantes equacions com paràmetres hi hagi i si fa falta utilitzem \textbf{la condició de derivabilitat} per treure una altra equació $\therefore$, derivada ha de ser contínua i les derivades laterals han de ser iguals.


\end{tcolorbox}
\subsection{Bolzano}
\begin{tcolorbox}[
  title=Teorema de Bolzano,
  colback=Green!5!white,
  colframe=Green!70!black,
  coltitle=white,
  colbacktitle=Green!70!black,
  fonttitle=\bfseries,
  arc=6pt,
  boxrule=0.8pt
]

El teorema de \textbf{Bolzano} garantitza que si en un interval de $[a,b]$ passa que $f(a)\cdot f(b) < 0$. Existeix almenys un punt $c \in [a,b]$ de manera que $f(c) =0$ per tant prenem 2 valors i anem \textbf{escurçant el interval}.   

\end{tcolorbox}
\section{T5 Identitats Trigonomètriques}
\subsection{Igualtats trigonomètriques}
\begin{tcolorbox}[
  title=Relacions trigonomètriques,
  colback=magenta!5!white,
  colframe=magenta!70!black,
  coltitle=white,
  colbacktitle=magenta!70!black,
  fonttitle=\bfseries,
  arc=6pt,
  boxrule=0.8pt
]

Podem relacionar les diferentes f(x) trigonomètriques com:

\[
\tan \theta = \frac{\sin \theta}{\cos \theta}
\]

\[
\sin^2 \theta + \cos^2 \theta = 1
\]

\[
1 + \tan^2 \theta = \sec^2 \theta
\]
\end{tcolorbox}
\subsection{Suma i resta d'angles}
\begin{tcolorbox}[
  title=Funcions trigonomètriques no perfectes,
  colback=violet!5!white,
  colframe=violet!70!black,
  coltitle=white,
  colbacktitle=violet!70!black,
  fonttitle=\bfseries,
  arc=6pt,
  boxrule=0.8pt
]

Si necessitem calcular $\sin \theta$, $\cos \theta$ i $\tan \theta$, on $\theta$ és una \textbf{combinació} de \textbf{dos angles}, ja sigui $\theta = \alpha + \beta$ o $\theta = 2\alpha$, utilitzem les identitats \textbf{trigonomètriques corresponents}.

\[
\sin(\alpha + \beta) = \sin \alpha \cos \beta \pm \cos \alpha \sin \beta
\]

\[
\cos(\alpha + \beta) = \cos \alpha \cos \beta \mp \sin \alpha \sin \beta
\]

\[
\sin(2\alpha) = 2 \sin \alpha \cos \alpha
\]

\[
\cos(2\alpha) = \cos^2 \alpha - \sin^2 \alpha
\]

\[
\sin^2(\alpha) = \frac{2 - \cos(2\alpha)}{2}
\]

\begin{center}
     $\uparrow$ (Super important per $\int$) $\uparrow$
\end{center}
\end{tcolorbox}

\section{T6-T7 Derivades i aplicacions a les derivadades}
\subsection{Derivades importants}
\begin{tcolorbox}[
  title=Derivades trivials,
  colback=orange!5!white,
  colframe=orange!70!black,
  coltitle=white,
  colbacktitle=orange!70!black,
  fonttitle=\bfseries,
  arc=6pt,
  boxrule=0.8pt
]

Les derivades que s'haurien de recordar, \textbf{a part de les bàsiques}, són les seguents:



\[
f(x) = \arccos(x) \therefore f'(x) = -\frac{1}{\sqrt{1-x^2}}
\]

\[
f(x) = \arcsin(x) \therefore f'(x) = \frac{1}{\sqrt{1-x^2}}
\]

\[
f(x) = \arctan(x) \therefore f'(x) = \frac{1}{1+x^2}
\]

\[
f(x) = \log_nx \therefore f'(x) = \frac{1}{x\ln(n)}
\]

\end{tcolorbox}
\subsection{Recta tangent}
\begin{tcolorbox}[
  title=Trobar equació de la recta tangent,
  colback=Peach!5!white,
  colframe=Peach!70!black,
  coltitle=white,
  colbacktitle=Peach!70!black,
  fonttitle=\bfseries,
  arc=6pt,
  boxrule=0.8pt
]

Trobar la \textbf{recta tangent} de la funció en un punt és molt fàcil, recodem \textit{$y=mx+n$} on \textbf{$m=f'(x_0)$} i fàcilment podem subsituir un punt \textbf{$P(x_0,f(x_0))$} $\in$ a la funció \textit{(el punt que ens demanen)} i aixi trobem $n$.

\end{tcolorbox}
\subsection{Optimització de funcions}
\begin{tcolorbox}[
  title=Problemes d'optimització,
  colback=BurntOrange!5!white,
  colframe=BurntOrange!70!black,
  coltitle=white,
  colbacktitle=BurntOrange!70!black,
  fonttitle=\bfseries,
  arc=6pt,
  boxrule=0.8pt
]

Aquests problemes són dificils i ens hem d'enrrecordar de les formules de \textbf{l'àrea, volum i superficie} de les figures geomètriques:\\
\\
\begin{center}
\includegraphics[width=0.5\textwidth]{Areas_Superficies.png}
\end{center}

A més ens poden demanar el Benefici com $B(x)$ que bàsicament és Ingrès - Costos $\therefore I(x) - C(x)$.
\\
Tot i sabent això el més important és \textbf{establir una relació entre les incògnites} per tenir una \textbf{funció en una variable}$\therefore$: \\
\\
1) Escriure la funció objectiu \textit{(és igual en quantes variables)}\\
2) Intentar \textbf{en el enunciat} trobar una relació ja sigui una igualtat o una relació geomètrica\\
3) Substituir a la funció f(x) objectiu\\
4) Trobar $f'(x)$ i $f'(x) = 0$ trobant els valors màxims i minims. ${\textbf{\textcolor{red}{Comprovar que és màxim o mínim.}}}$\\

Tips: \\
· Sobre tot \textbf{dibuixar la situació}.\\
· Si tenim una recta busquem un punt $P(x,f(x))$, buscant \textbf{l'equació de la recta.}\\
· Si hem de buscar una \textbf{distància punt recta}, usem la formula $\Rightarrow d(P,r) = \frac{\lvert Ax_0 + By_0 + C \rvert}{\sqrt{A^2 + B^2}} $, posant el punt en $P(x_0,y_0)$ i l'equació en forma general $Ax + By +C = 0$. 

\end{tcolorbox}

\subsection{Punts d'estudi}

\begin{tcolorbox}[
  title=Punts d'inflexió,
  colback=Apricot!5!white,
  colframe=Apricot!70!black,
  coltitle=white,
  colbacktitle=Apricot!70!black,
  fonttitle=\bfseries,
  arc=6pt,
  boxrule=0.8pt
]

Si $f'(x) = 0$ por ser \textbf{màxim, mínim o p.inflexió}, podem estudiar aquests casos amb $f''(x)$ que ens descriu \textbf{la curvatura}.\\
\\
Si $f'(x)=0$ també hem d'estudiar $f''(x) = 0$ ja \textbf{que potser no comparteixen }punts en comú que $\exists$, també escriure els intèrvals en forma $(-\infty , x) \cup (y, +\infty)$, {\textbf{\textcolor{red}{L'infinit sempre va amb ().}}}\\
\\
1) Si $f''(x) > 0 \iff$ tenim un mínim. (\textbf{Convexa})\\
\begin{center}
\begin{tikzpicture}
\draw[->] (-2,0) -- (2,0) node[right] {$x$};
\draw[->] (0,-1) -- (0,3) node[above] {$y$};
\draw[thick, color = orange, domain=-1.5:1.5, samples=100] plot (\x,{\x*\x});
\node at (1.5,2.5) {$f''(x)>0$};
\end{tikzpicture}\\
\end{center}
2) Si $f''(x) < 0 \iff$ tenim un màxim.(\textbf{Còncava})\\
\begin{center}
\begin{tikzpicture}
\draw[->] (-2,0) -- (2,0) node[right] {$x$};
\draw[->] (0,-3) -- (0,1) node[above] {$y$};
\draw[thick, color = orange, domain=-1.5:1.5, samples=100] plot (\x,{-(\x*\x)});
\node at (1.5,-2.5) {$f''(x)<0$};
\end{tikzpicture}
\end{center}
3) Si $f''(x) = 0 \iff$ tenim un p.inflexió.\\
\begin{center}
\begin{tikzpicture}[scale=0.7]
  \begin{axis}[
    axis lines = middle,
    xlabel=$x$, ylabel=$y$,
    domain=-2:2, samples=200,
    title={$f(x)=x^3$}
  ]
  \addplot[thick, color = orange] {x^3};
  \addplot[only marks, mark=*] coordinates {(0,0)};
  \end{axis}
\end{tikzpicture}
\end{center}
\end{tcolorbox}

\section{T8 Probabilitat}
\subsection{Formulari de probabilitat}

\begin{tcolorbox}[
  title=Formules Probabilitat,
  colback=Goldenrod!5!white,
  colframe=Goldenrod!70!black,
  coltitle=white,
  colbacktitle=Goldenrod!70!black,
  fonttitle=\bfseries,
  arc=6pt,
  boxrule=0.8pt
]

A la probabilitat de la selectivitat tenim diferents formules que ens puguin sortir. Primer el que s'ha de fer es \textbf{definir l'esdeveniment} i apuntar totes les probabilitats que ens diu l'anunciat i els seus \textit{complementaris}. $P(A^c) = 1- P(A)$ i $P(B^c\mid A) = 1-P(B\mid A)$.\\
\\
Les formules que ens hem de saber són \textbf{Probabilitat total} i \textbf{Bayes};

Mentre que la \textit{probabilitat total} ens serveix per \textbf{extreure la probabilitat d'una condicionada}, \textit{Bayes} ens serveix perquè es giri\textbf{ l'ordre de dependència.}

\[
P(A) = \frac{\mid A \mid}{\mid \Omega \mid} 
\]

\[
P(B) = P(A)·P(B\mid A) + P(A^c)·P(B\mid A^c)...
\]

\[
P(B\mid A) = \frac{P(B)·P(A\mid B)}{P(A)}
\]

\[
P(A\mid B) = \frac{P(A\cap B)}{P(B)}
\]

\[
P(A\cup B) = P(A)+P(B)-P(A\cap B)
\]
\end{tcolorbox}

\section{T9-T10-T11 Geometria}
\subsection{Usos dels vectors}

\begin{tcolorbox}[
  title=Utlitats dels vectors,
  colback=Brown!5!white,
  colframe=Brown!70!black,
  coltitle=white,
  colbacktitle=Brown!70!black,
  fonttitle=\bfseries,
  arc=6pt,
  boxrule=0.8pt
]

\textbf{Punt mitjà}\\
\\
Calculem el \textbf{punt mitjà}; $\vec{OP} = \vec{OA}+\frac{1}{2}\vec{AB} = P \iff P = \frac{A+B}{2}$\\
\\
\textbf{Proporcionalitat}\\
\\
Dos vectors son \textbf{paralels}, si son proporcionals; $\frac{\vec{v_x}}{\vec{u_x}} = \frac{\vec{v_y}}{\vec{u_y}} = \frac{\vec{v_z}}{\vec{u_z}} = k \iff k\in \mathbb{R}$
\\

\end{tcolorbox}

\subsection{Vectors perpendiculars (base)}

\begin{tcolorbox}[
  title=Linealitat de vectors,
  colback=Maroon!5!white,
  colframe=Maroon!70!black,
  coltitle=white,
  colbacktitle=Maroon!70!black,
  fonttitle=\bfseries,
  arc=6pt,
  boxrule=0.8pt
]

\noindent En $R^2$, dos vectors \textbf{son l. independents} si no es poden formar l'un per l'altre és a dir $\therefore$ $\vec{v} = k\vec{b} \longleftrightarrow k\in \mathbb{R}$, per tant si dividim $\frac{\vec{v}}{\vec{b}} = k$ per totes les components.\\

\noindent En $R^3$, necessitem 3 vectors per formar una \textbf{base}, si $\det(V) \neq 0$, si V = matriu de vectors (són linealment independents). Podem trobar si son l.d per que estan formats fent un sistema $\vec{v} = x\vec{u}+y\vec{w}$

\end{tcolorbox}

\subsection{Canvis de base}

\begin{tcolorbox}[
  title=Com fer canvis de Bases,
  colback=Sepia!5!white,
  colframe=Sepia!70!black,
  coltitle=white,
  colbacktitle=Sepia!70!black,
  fonttitle=\bfseries,
  arc=6pt,
  boxrule=0.8pt
]

Si volem \textbf{expressar un vector} $\vec{a}=(1,2,-3)$ en una base $B =\{(1,-1,3),(1,0,0), (2,1,0)\}$\\
\\
$\vec{a} =\lambda\vec{v}+k\vec{u}+\mu\vec{w} \iff$ resolem el sistema fins trobar $\lambda, k, \mu$, i operem per trobar $\vec{a}$\\
\\
Si tenim un vector $\vec{a}=(1,2,-3)$ en una base $B =\{(1,-1,3),(1,0,0), (2,1,0)\}$, el vector $\vec{a}'$ \textbf{en cordenades cartesianes} serà el resultant de; \\
\\
$\vec{a}' = a_x (1,-1,3) + a_y(1,0,0)+a_z(2,1,0)$ = $\vec{a}' = 1 (1,-1,3) + 2(1,0,0)+-3(2,1,0)$


\end{tcolorbox}

\subsection{Producte escalar}

\begin{tcolorbox}[
  title=Producte escalar i usos,
  colback=Brown!5!white,
  colframe=Brown!70!black,
  coltitle=white,
  colbacktitle=Brown!70!black,
  fonttitle=\bfseries,
  arc=6pt,
  boxrule=0.8pt
]

Anomenem producte escalar a $\vec{v}\cdot \vec{u} = \mid\vec{v}\mid\mid\vec{u}\mid\cos(\alpha)$, es pot utilitzar per; \\
\\
\textbf{Angle entre vectors}\\
\\
$\alpha = \arccos(\frac{\vec{v}\cdot \vec{u}}{\mid\vec{v}\mid\mid\vec{u}\mid}) \iff 0º  \leq \alpha \leq 180º $ \\
\\
\textbf{Projecció d'un vector sobre un altre}\\
\\
$Proj_{\vec{v}}\vec{u} = \mid\vec{u}\mid\cos(\alpha)\cdot\frac{\vec{v}}{\mid\vec{v}\mid}$\\
\\

\end{tcolorbox}

\subsection{Producte vectorial}

\begin{tcolorbox}[
  title=Producte vectorial i usos,
  colback=Mahogany!5!white,
  colframe=Mahogany!70!black,
  coltitle=white,
  colbacktitle=Mahogany!70!black,
  fonttitle=\bfseries,
  arc=6pt,
  boxrule=0.8pt
]

Anomenem producte vectorial a $\vec{v}\times \vec{u} =
=
\begin{vmatrix}
\mathbf{i} & \mathbf{j} & \mathbf{k} \\
v_x & v_y & v_z \\
u_x & u_y & u_z
\end{vmatrix}
$, es pot utilitzar per;\\
\\
\textbf{Càlcul del paralelogram}\\
\\
$Area_{paralelogram} = \mid \vec{v}\times \vec{u}\mid$ \\
$Area_{triangle} = \frac{1}{2}\mid \vec{v}\times \vec{u}\mid$ \\
\\
\textbf{Càlcul d'una base}\\
\\
Si volem trobar una base i ens donen un vector $\vec{v} = (v_x,v_y,v_z)$, trobem un vector $\therefore \vec{v}\cdot\vec{u} = 0 \iff \vec{u} = (v_y,-v_x,0)$, i per acabar trobem un vector fent el \textbf{producte vectorial d'aquests dos}.
\\

\end{tcolorbox}

\subsection{Producte mixt}

\begin{tcolorbox}[
  title=Producte mixt i usos,
  colback=Sienna!5!white,
  colframe=Sienna!70!black,
  coltitle=white,
  colbacktitle=Sienna!70!black,
  fonttitle=\bfseries,
  arc=6pt,
  boxrule=0.8pt
]

Anomenem producte mixt a una operació entre 3 vectors; $\vec{u}, \vec{v}, \vec{w}$ i s'escriu tal que $[\vec{u}, \vec{v}, \vec{w}] = \vec{u}\cdot(\vec{v}\times\vec{w})$, aquest producte no commuta i bàsicament ens decriu geometricament \textbf{el volum del paral·lelopìptede} que formen els 3 vectors $\therefore$ (les unitats solen ser $u^2$).

\begin{center}
\includegraphics[width=0.3\textwidth]{Producte_Mixt.png}
\end{center}

Una forma més fàcil de calcular-lo és fent el $ \det(V) \longleftrightarrow V = \begin{vmatrix}
v_1 \\
v_2 \\
v_3 \\
\end{vmatrix}$, pot $\exists$ una àrea negativa, que \textbf{geometricament} vol dir que el vector ($\vec{u}$), està en sentit contrari al vector resultant ($\vec{v}\times\vec{w}$)

\end{tcolorbox}

Cal saber que l'area del tetraedre és composta per $\frac{1}{6}\lvert [\vec{v},\vec{u},\vec{w}] \rvert$, és dir que té 2 solucions ja que \[
|x|=\begin{cases}
x & \text{si } x \ge 0 \\
-x & \text{si } x < 0
\end{cases}
\]

\subsection{Equacions de la recta}

\begin{tcolorbox}[
  title=Equacions de la recta,
  colback=SaddleBrown!5!white,
  colframe=SaddleBrown!70!black,
  coltitle=white,
  colbacktitle=SaddleBrown!70!black,
  fonttitle=\bfseries,
  arc=6pt,
  boxrule=0.8pt
]

Podem definir qualsavol \textbf{recta}, amb (\textbf{Un punt i un vector director}),(\textbf{Dos punts}), les rectes tenen infinits punts al espai $\therefore$ hi ha infinites cominacions d'equacions

\begin{center}
\includegraphics[width=0.3\textwidth]{Equacions_d'una_recta.png}
\end{center}

\begin{center}
\renewcommand{\arraystretch}{1.6}
\begin{tabular}{|c|c|}
\hline
\textbf{\textit{Equació}} & \textbf{\textit{Nom}} \\
\hline

$r: (x,y,z) = (x_0,y_0,z_0) + \lambda (v_x,v_y,v_z)$ 
& Equació vectorial \\

\hline

$\begin{cases}
x = x_0 + \lambda v_x \\
y = y_0 + \lambda v_y \\
z = z_0 + \lambda v_z
\end{cases}$
& Equacions paramètriques \\

\hline

$\dfrac{x - x_0}{v_x}
=
\dfrac{y - y_0}{v_y}
=
\dfrac{z - z_0}{v_z}$
& Equació contínua \\

\hline

$\begin{cases}
   A x + B y + C = 0 \\
A' x + B' y + C' = 0
\end{cases}$

& Equació general (implícita) \\

\hline

\end{tabular}
\end{center}


\end{tcolorbox}

\subsection{Equacions d'un pla}

\begin{tcolorbox}[
  title=Equacions d'un pla,
  colback=DarkRed!5!white,
  colframe=DarkRed!70!black,
  coltitle=white,
  colbacktitle=DarkRed!70!black,
  fonttitle=\bfseries,
  arc=6pt,
  boxrule=0.8pt
]

Definim un pla de dues maneres: (\textbf{Un punt i dos vectors l.i.}, \textbf{3 punts no alineats}). El definim de forma:

\begin{center}
\includegraphics[width=0.3\textwidth]{Equacions_del_Pla.png}
\end{center}

\[
\pi : (x,y,z) = (x_0,y_0,z_0) + 
\lambda (v_x, v_y, v_z) + 
\mu (u_x, u_y, u_z)
\]

o com un \textbf{sistema}.

\[\begin{cases}
    x = x_0 + \lambda v_x + \mu u_x \\
    y = y_0 + \lambda v_y + \mu u_y \\
    z = z_0 + \lambda v_z + \mu u_z
\end{cases}\]\\

Finalment el podem expressar amb la forma general de \textbf{$Ax+By+Cz+D =0$} on $(A,B,C)$ és el vector resultant del \textbf{producte vectorial entre els vectors directos del pla}. i D és una variable que trobem substituint x,y,z per un punt que pertany al pla.

\end{tcolorbox}

\subsection{Projeccions}

\begin{tcolorbox}[
  title= Estudi de projeccions ortogonals,
  colback=Mahogany!5!white,
  colframe=Mahogany!70!black,
  coltitle=white,
  colbacktitle=Mahogany!70!black,
  fonttitle=\bfseries,
  arc=6pt,
  boxrule=0.8pt
]

\textbf{Punt sobre recta:}
\begin{enumerate}
    \item Creem un pla perpendicular a r que passi pel punt
    \item Fem la intersecció del pla i la recta ($\pi \cap s$)
\end{enumerate}
\textbf{Punt sobre pla:}
\begin{enumerate}
    \item Creem una recta que amb el vector normal i el punt
    \item Busquem la intersecció de la recta i el pla ($\pi \cap s$)
\end{enumerate}
\textbf{Recta sobre pla}
\begin{enumerate}
    \item Fem la Projecció de dos punts sobre el pla
    \item Creem una recta amb aquests dos punts
\end{enumerate}
\end{tcolorbox}

\subsection{Posicions relatives}

\begin{tcolorbox}[
  title=Estudi de posicions,
  colback=Bittersweet!5!white,
  colframe=Bittersweet!70!black,
  coltitle=white,
  colbacktitle=Bittersweet!70!black,
  fonttitle=\bfseries,
  arc=6pt,
  boxrule=0.8pt
]

\textbf{Dos rectes:}

\begin{center}
\begin{tabular}{lcr}
\noalign{\hrule height 0.8pt}
\textbf{Posició} & \textbf{Condició} \\
\noalign{\hrule height 0.3pt}
\\[-2mm] % reduce space if needed
Paral·leles & $\vec{v_1} \parallel \vec{v_2}$ i $P_1 \notin r_2$ \\
Coincidents & $\vec{v_1} \parallel \vec{v_2}$ i $P_1 \in r_2$ \\
Secants & $\vec{v_1} \nparallel \vec{v_2}$ i $[\vec{PQ},\vec{v_1},\vec{v_2}] = 0$ \\
Encreuades & $\vec{v_1} \nparallel \vec{v_2}$ i $[\vec{PQ},\vec{v_1},\vec{v_2}] \neq  0$\\
\noalign{\hrule height 0.8pt}
\end{tabular}
\end{center}

\textbf{Dos plans:}

\begin{center}
\begin{tabular}{lcr}
\noalign{\hrule height 0.8pt}
\textbf{Posició} & \textbf{Condició} \\
\noalign{\hrule height 0.3pt}
\\[-2mm] % reduce space if needed
Paral·lels & $\vec{n_1} \parallel \vec{n_2}$ i $P_1 \notin \pi_2$ \\
Coincidents & $\vec{n_1} \parallel \vec{n_2}$ i $P_1 \in \pi_2$ \\
Secants & $\vec{n_1} \nparallel \vec{n_2}$ i tenen una recta en comú \\
\noalign{\hrule height 0.8pt}
\end{tabular}
\end{center}

\textbf{Una recta i un pla:}

\begin{center}
\begin{tabular}{lcr}
\noalign{\hrule height 0.8pt}
\textbf{Posició} & \textbf{Condició} \\
\noalign{\hrule height 0.3pt}
\\[-2mm] % reduce space if needed
Paral·lels & $\vec{n} \cdot \vec{v} = 0$ i $P \notin \pi$ \\
Coincidents & $\vec{n} \cdot \vec{v} = 0$ i $P \in \pi$ \\
Secants & $\vec{n} \cdot \vec{v} \neq 0$ \\
\noalign{\hrule height 0.8pt}
\end{tabular}
\end{center}

\textbf{Tres Plans:} Resolem el sistema dels tres plans;
\\
Siguin $\pi_1,\pi_2,\pi_3$ tres plans a $\mathbb{R}^3$.

\begin{itemize}
\item[$\star$] \textbf{SDC}

\item[$\cdot$] Una única solució. $P(x_0,y_0,z_0)$

\item[$\star$] \textbf{SCI}

\item[$\cdot$] 3 plans coincidents. $\frac{A_1}{A_2} = \frac{B_1}{B_2} = \frac{C_1}{C_2} = \frac{D_1}{D_2}$ i $\frac{A_1}{A_2} = \frac{B_1}{B_3} = \frac{C_1}{C_3} = \frac{D_1}{D_3}$ i $\vec{n_1} \parallel \vec{n_2} \parallel \vec{n_3}$
\item[$\cdot$] 2 coincidents i 1 secant. $\frac{A_1}{A_2} = \frac{B_1}{B_2} = \frac{C_1}{C_2} = \frac{D_1}{D_2}$ i $\vec{n_1} \parallel \vec{n_2} \nparallel \vec{n_3}$
\item[$\cdot$] Tallen en una recta. $[\vec{n_1},\vec{n_2},\vec{n_3}]$

\item[$\star$] \textbf{SI}

\item[$\cdot$] 3 plans paral·lels. $\frac{A_1}{A_2} = \frac{B_1}{B_2} = \frac{C_1}{C_2} \neq \frac{D_1}{D_2}$ i $\frac{A_1}{A_2} = \frac{B_1}{B_3} = \frac{C_1}{C_3} \neq \frac{D_1}{D_3}$
\item[$\cdot$] 2 coincidents i 1 paral·lel. $\frac{A_1}{A_2} = \frac{B_1}{B_2} = \frac{C_1}{C_2} = \frac{D_1}{D_2}$ i $\frac{A_1}{A_2} = \frac{B_1}{B_3} = \frac{C_1}{C_3} \neq \frac{D_1}{D_3}$
\item[$\cdot$] 2 paral·lels i 1 secant. $\vec{n_1} \parallel \vec{n_2} \nparallel \vec{n_3}$
\item[$\cdot$] Secants 2:2. $[\vec{n_1}, \vec{n_2}, \vec{n_3}] = 0$

\end{itemize}

\end{tcolorbox}

\subsection{Distàncies}

\begin{tcolorbox}[
  title=Producte escalar i usos,
  colback=RawSienna!5!white,
  colframe=RawSienna!70!black,
  coltitle=white,
  colbacktitle=RawSienna!70!black,
  fonttitle=\bfseries,
  arc=6pt,
  boxrule=0.8pt
]

Podem trobar diferentes distàncies, ens podem aprendre les formules o fer \textbf{projecció i el mòdul del vector resultant}, les formules són:\\

\textbf{Distància Punt-Punt:}\\
\\
$d(P, Q) = \sqrt{(x_q-x_p)^2+(y_q-y_p)^2(z_q-z_p)^2}$\\
\\
\textbf{Distància Punt-Recta:}\\
\\
$d(P,r) = \frac{\lvert \vec{PR} \times \vec{v_r} \rvert}{\lvert \vec{v_r}\rvert}$\\
\\
\textbf{Distància Recta-Recta:}\\
\\
$d(r,s) = \frac{\lvert \vec{v_r} \times \vec{P_rP_s}\rvert}{\lvert \vec{v_r}\rvert}$\\
\\
\textbf{Distància Recta-Recta paral·leles:}\\
\\
$d(r,s) = \frac{\lvert [\vec{v_r},\vec{v_s}, \vec{P_rP_s}]\rvert}{\lvert\vec{v_r}\times\vec{v_s}\rvert}$\\
\\
\textbf{Distància Punt-Pla:}\\
\\
$d(P,\pi) = \frac{\lvert Ax_p+By_p+Cz_p+D \rvert}{\sqrt{A^2+B^2+C^2}}$\\
\\
\end{tcolorbox}

\section{T12 Integrals}
\subsection{Integrals importants}
\begin{tcolorbox}[
  title=Integrals trivials,
  colback=Orchid!5!white,
  colframe=Orchid!70!black,
  coltitle=white,
  colbacktitle=Orchid!70!black,
  fonttitle=\bfseries,
  arc=6pt,
  boxrule=0.8pt
]

Les integrals que s'haurien de recordar, \textbf{a part de les bàsiques}, són les seguents:

\[
\int \frac{1}{\cos ^2x} dx = \tan x + K
\]

\[
\int \frac{1}{\sqrt{1-x^2}} dx = \arcsin x + K
\]

\[
\int -\frac{1}{\sqrt{1-x^2}} dx = \arccos x + K
\]

\[
\int \frac{1}{1+x^2} dx = \arctan x + K
\]

Sempre posar "+K" per a \textbf{qualsevol constant al fer la operació inversa}, i que $K \in \mathbb{R}$

\end{tcolorbox}

\subsection{Tècniques d'integració}

\begin{tcolorbox}[
  title=Com integrar,
  colback=CarnationPink!5!white,
  colframe=CarnationPink!70!black,
  coltitle=white,
  colbacktitle=CarnationPink!70!black,
  fonttitle=\bfseries,
  arc=6pt,
  boxrule=0.8pt
]

Les integrals que s'haurien de recordar, \textbf{a part de les bàsiques}, són les seguents:

\begin{itemize}
    \item \textbf{U-Sub}, $f(x) = u, du = f'(x) dx$ i aïllem dx, substituïm u a la integral i integrem.

    \item \textbf{Integrals per parts}, Usar la formula $\int u\cdot dv = u\cdot v \int v\cdot du$, intentar reduïr el polinomi (si n'hi ha a 0) i vigilar amb les cicliques

    \item \textbf{Fraccions parcials}, Tenim diversos cassos (vigilar perque potser tenim una polinòmica o quasi-immmediata); $\int\frac{p(x)}{q(x)}$

    \begin{itemize}
        \item Grau num $<$ Grau denom $\int \frac{x_1^{n-1} + x_2^{n-2} ...+x_m^0}{x_1^{n} + x_2^{n-1} ...+x_m^0}$ 
        
        \begin{itemize}
            \item Segurament es un polinomi, un $ln\abs{x}$ o una $\arctan(x)$
            \item Fer la tècnica de $x-1 = Ax+B(x+2)$, \textbf{donem valors a x i aïllem A i B}.
                \begin{itemize}
                    \item Si son repetides, anem de $x .. x^n$ i afegim una fracció per cada
                    \item Si hi ha una arrel no real ($q(x)$ en forma $x^2+a^2$) dividirem la fracció amb les reals com A i B i les no reals $\frac{Mx+N}{x^2+a^2}$ i ara podem derivar-la amb $ln\abs{x}$ o $\arctan(x)$
                \end{itemize}
        \end{itemize} 
        \item Grau num $\geq$ Grau denom $\int \frac{x_1^{n} + x_2^{n-1} ...+x_m^0}{x_1^{n-1} + x_2^{n-2} ...+x_m^0}$, podem reescriure i tornar a integrar el numerador com a ($p(x) = d(x)\cdot q(x) +r(x)$) on d(x) és el resultat i q(x) el divisor
        \begin{itemize}
            \item Fem divisió per \textbf{ruffini}, si tenim la foram $(x\pm a)$
            \item Fem divisió per caixeta
        \end{itemize}

    \end{itemize}

    
\end{itemize}
\end{tcolorbox}

\subsection{Integrals definides}

\begin{tcolorbox}[
  title=Llei de Barrow,
  colback=Purple!5!white,
  colframe=Purple!70!black,
  coltitle=white,
  colbacktitle=Purple!70!black,
  fonttitle=\bfseries,
  arc=6pt,
  boxrule=0.8pt
]

\textbf{Barrow}, ens diu que $\int_a^b f(x) dx = F(b)-F(a)$,  {\textbf{\textcolor{red}{tenir en ment que $\exists$ A $\leq$ 0. Implica que esta sota l'eix.}}}\\

\begin{itemize}
    \item $\int_{a}^{a} f(x)  dx = 0$
     \item $\int_{a}^{a} f(x)  dx = 0$
     \item Si $c\in (a,b)$ $\int_a^b f(x) dx = \int_a^c f(x) dx + \int_c^b f(x) dx$
     \item $\int_{a}^{b} f(x)  dx = -\int_{b}^{a} f(x)  dx $\\
\end{itemize}

Per calcular l'area compressa entre \textbf{dos corbes}, mirem els punts de tall com a \textit{limits d'integració}.
Fem la integral de $\int_a^b [f(x)-g(x)] dx$ on f(x) es superior a g(x).\\

Per calcular l'area compresa entre \textbf{tres funcions}, prenem per intèrvals les \textit{dos funcions més altes} i fem l'area compressa entre les dues $\therefore A_t = \sum_{i=1}^{n}
A_i$.

\vspace{0.6cm}

\definecolor{lilaoscuro}{HTML}{5B2D8E}
\definecolor{lilamedio}{HTML}{9B59B6}
\definecolor{lilaclaro}{HTML}{C39BD3}
\definecolor{lilafondo}{HTML}{F3E5F5}
\definecolor{lilarelleno}{HTML}{CE93D8}
\definecolor{lilatexto}{HTML}{4A235A}
\definecolor{fondopag}{HTML}{FDFAFF}

\begin{center}
\begin{tikzpicture}

\begin{axis}[
  width  = 9cm,
  height = 9cm,
  xmin = -0.25, xmax = 1.35,
  ymin = -0.15, ymax = 1.25,
  axis lines = center,
  axis line style = {color=lilaoscuro, line width=1.2pt},
  tick style     = {color=lilaoscuro, line width=0.8pt},
  ticklabel style= {color=lilatexto, font=\small},
  xtick = {0, 0.25, 0.5, 0.75, 1},
  ytick = {0, 0.25, 0.5, 0.75, 1},
  xlabel = {$x$},
  ylabel = {$y$},
  xlabel style = {color=lilaoscuro, font=\bfseries\large,
                  at={(axis cs:1.35,0)}, anchor=north},
  ylabel style = {color=lilaoscuro, font=\bfseries\large,
                  at={(axis cs:0,1.25)}, anchor=east},
  grid = both,
  grid style = {color=lilarelleno!35, line width=0.4pt},
  major grid style={color=lilarelleno!55, line width=0.6pt},
  minor tick num = 1,
  legend pos = north west,
  legend style = {
    fill=lilafondo,
    draw=lilaoscuro!60,
    rounded corners=3pt,
    font=\small,
    row sep=3pt
  },
  clip = false,
]

\addplot [name path=F, draw=none, domain=0:1, samples=120] {x};
\addplot [name path=G, draw=none, domain=0:1, samples=120] {x^2};

\addplot [
  fill=lilarelleno,
  fill opacity=0.45,
] fill between [of=F and G];

\addplot [
  domain = -0.1:1.05,
  samples = 200,
  color = lilamedio,
  line width = 2.2pt,
] {x^2};

\addplot [
  domain = -0.1:1.05,
  samples = 200,
  color = lilaclaro!80!lilaoscuro,
  line width = 2.2pt,
  dashed,
] {x};

\addplot [
  only marks,
  mark = *,
  mark size = 3.8pt,
  color = lilaoscuro,
] coordinates {(0,0) (1,1)};

\node[above left, color=lilaoscuro, font=\small\bfseries]
  at (axis cs:0,0) {$(0,0)$};

\node[above left, color=lilaoscuro, font=\small\bfseries]
  at (axis cs:1,1) {$(1,1)$};


\addplot [
  color=lilaoscuro!50,
  line width=0.9pt,
  dotted,
] coordinates {(1,0) (1,1)};


\end{axis}
\end{tikzpicture}
\end{center}

\vspace{0.5cm}

\end{tcolorbox}
\end{document}